\documentclass[11pt,a4paper]{article}
\usepackage[utf8]{inputenc}
\usepackage[T1]{fontenc}
\usepackage[margin=0.65in]{geometry}
\usepackage{enumitem}
\usepackage{titlesec}
\usepackage{xcolor}
\usepackage{hyperref}
\usepackage{helvet}
\renewcommand{\familydefault}{\sfdefault}

\definecolor{primarycolor}{RGB}{33,37,41}
\definecolor{secondarycolor}{RGB}{0,92,150}
\definecolor{accentcolor}{RGB}{90,100,110}
\definecolor{sectioncolor}{RGB}{0,80,40}
\hypersetup{colorlinks=true,urlcolor=secondarycolor}
\titleformat{\section}{\large\bfseries\color{sectioncolor}}{}{0em}{}
\titlespacing*{\section}{0pt}{0.65em}{0.25em}
\setlist[itemize]{leftmargin=1.2em,itemsep=0.1em,topsep=0.15em}
\pagestyle{empty}

\begin{document}

\begin{center}
{\LARGE\bfseries\color{primarycolor} MARTA KOWAL-BRANDT}\\[0.25em]
{\large QA Automation Engineer | CI/CD and Test Infrastructure}\\[0.25em]
{\small Wroclaw, Poland | marta.kowalbrandt@example.com | +48 555 303 303 | gitlab.com/mkb-example}
\end{center}
\vspace{-0.4em}

\section*{PROFESSIONAL SUMMARY}
QA automation engineer with experience owning CI/CD pipelines and building test infrastructure for three product teams. Reduced average pipeline time from 38 to 14 minutes and improved end-to-end test reliability from a 7\% flaky rate to below 1\%. Hands-on with GitLab CI, Docker, Helm and namespace-level Kubernetes deployments, with additional personal experience operating a k3s GitOps homelab and provisioning a Hetzner VPS with Terraform. Seeking a move into a DevOps or platform engineering role.

\section*{TECHNICAL SKILLS}
\textbf{CI/CD and containers:} GitLab CI, Jenkins, Docker, Docker Compose\\
\textbf{Infrastructure:} Helm, Kubernetes (namespace-level access and test-environment deployments), k3s, Terraform (personal project)\\
\textbf{Monitoring and TLS:} Prometheus, Grafana, Loki, cert-manager, TLS certificate management\\
\textbf{Cloud learning:} AWS EC2 and S3 fundamentals (personal learning)\\
\textbf{Programming and data:} Python, Bash, Java (test code), SQL, PostgreSQL\\
\textbf{Testing:} pytest, Playwright, Selenium, Postman\\
\textbf{Other:} Linux, Git, Ansible (basic)

\section*{PROFESSIONAL EXPERIENCE}
\textbf{QA Automation Engineer (Test Infrastructure) | Odrava Health Software, Wroclaw | 2021 - Present}
\begin{itemize}
\item Own GitLab CI pipelines for three product teams (approximately 25 developers), covering build, unit, integration and end-to-end test stages and artifact publishing.
\item Built short-lived merge-request environments for the application, PostgreSQL and seeded data, deployed via Helm to the shared Kubernetes cluster from GitLab CI; environments are removed automatically after merge or 48 hours of inactivity.
\item Reduced average pipeline time from 38 to 14 minutes by profiling pipelines, parallelising test jobs, caching dependencies and separating end-to-end tests; reduced monthly CI runner costs by about 30\%.
\item Maintain a pytest and Playwright end-to-end suite of approximately 900 tests; reduced the flaky rate from 7\% to below 1\%.
\item Wrote Python and Bash tooling for test-data seeding; collected CI test results in PostgreSQL and built Grafana dashboards to track flaky tests.
\item Wrote internal documentation on how CI works, used in onboarding, and mentored two QA interns in 2023.
\end{itemize}

\textbf{QA Engineer | Brightmill Software House, Wroclaw | 2019 - 2021}
\begin{itemize}
\item Tested web and mobile applications for client projects using manual and automated testing.
\item Introduced a Java and Selenium automated regression suite for the largest client project.
\item Configured Jenkins jobs for nightly regression runs.
\end{itemize}

\section*{PROJECTS}
\textbf{Homelab k3s Cluster with GitOps | Personal project | 2024 - 2025}
\begin{itemize}
\item Ran a three-node Raspberry Pi 5 k3s cluster, deploying applications from a Git repository with Argo CD.
\item Set up Prometheus, Grafana and Loki monitoring, Traefik ingress and cert-manager for TLS; provisioned a small Hetzner VPS with Terraform and used Ansible (basic).
\end{itemize}

\section*{EDUCATION}
BEng Computer Engineering | Silesian Lowlands University of Technology | 2018

\section*{CERTIFICATION AND PROFESSIONAL DEVELOPMENT}
ISTQB Foundation Level | 2019\\
Preparing for Certified Kubernetes Administrator (CKA); exam scheduled for January 2027 (not yet passed).

\section*{LANGUAGES}
Polish (native), English (C1), German (A2)

\end{document}